\begin{resume}
% 明评完整显示，同参考文献的写法,在成果最后可附上SCI、EI检索号、影响因子等标志成果学术水平。
  
  \section*{发表的学术论文} % 发表的和录用的合在一起

  \begin{enumerate}[label={[\arabic*]},itemsep=0pt,parsep=0pt,labelindent=26pt,labelwidth=*,leftmargin=0pt,itemindent=*,align=left]

  \item Dan Wu, Kui Dai, Zhiying Wang. Retargetable Machine-Description System: Multi-layer Architecture Approach. 4th International Conference of on Grid and Cooperative Computing, November/December 2005, Beijing, China. LNCS 3795 Springer-Verlag, 1161~1166, ISSN: 0302-9743. (SCI Index, IDS N.O: BDQ17).
  \item 吴先宇,罗世彬,陈小前等.基于替代模型的高超声速进气道优化[J].弹箭与制导学报,2008,28(1):148-152.
  \item 第一作者.基于替代模型的高超声速进气道优化[J].弹箭与制导学报，2008,28(1):148-152.

  \end{enumerate}

  \section*{研究成果} % 有就写，没有就删除
  \begin{enumerate}[label={[\arabic*]},itemsep=0pt,parsep=0pt,labelindent=26pt,labelwidth=*,leftmargin=0pt,itemindent=*,align=left]

  \item 任天令, 杨轶, 朱一平, 等. 硅基铁电微声学传感器畴极化区域控制和电极连接的
    方法: 中国, CN1602118A. (中国专利公开号.)
  \item Ren T L, Yang Y, Zhu Y P, et al. Piezoelectric micro acoustic sensor
    based on ferroelectric materials: USA, No.11/215, 102. (美国发明专利申请号.)

  \end{enumerate}
\end{resume}
